%======================================================================
%  Focused experimental design
%======================================================================

\section{Focused experimental design}
\label{sec:methodology}

\subsection{Objective and research questions}

The thesis should make a narrower claim than ``CPUs can run large
language models''.  Its objective is to quantify when and why
on-package HBM changes inference on a Xeon Max processor, using DDR5 on
the same socket as the control.  This produces three primary research
questions:

\begin{table}[tbp]
\centering
\small
\caption{Primary research questions and predictions.  Predictions are
registered expectations to test, not conclusions.}
\label{tab:research-questions}
\begin{tabularx}{\textwidth}{@{}lYY@{}}
\toprule
 & \textbf{Question} & \textbf{Prediction} \\
\midrule
RQ1
 & How does verified local HBM change prefill and decode performance
   relative to local DDR5 on the same Max 9480?
 & The largest benefit occurs in low-batch decode.  Prefill and batched
   decode show less benefit as weight reuse and matrix-engine utilization
   increase. \\
RQ2
 & How do weight precision, KV precision and live-set size change the
   result as the working set approaches and exceeds 64\,GB per socket?
 & Quantization moves the capacity boundary but does not yield a
   bit-width-proportional speedup.  Crossing local HBM introduces a
   placement-dependent performance change. \\
RQ3
 & Can explicit socket or SNC-domain placement use additional cores and
   bandwidth without the regressions reported for unmanaged execution?
 & Local replicas or shards reduce UPI traffic and outperform a
   placement-oblivious two-socket run; SNC4 is expected to help when
   allocation and workers share the same domains. \\
\bottomrule
\end{tabularx}
\end{table}

Energy per token is a conditional outcome, not a primary research
question until a user-accessible measurement interface has been
validated.  Multi-node inference, GPU comparison and online serving are
outside the minimum scope.  They should be added only if the single-node
results create a specific need or hypothesis.

\FloatBarrier

\subsection{Evidence and the primary contrast}

Four evidence levels should remain distinct:

\begin{enumerate}
  \item a product guide establishes what a platform \emph{can} support;
  \item ENEAGRID documentation establishes what the site reported at a
        stated date;
  \item the system manifest establishes what an allocated node exposed;
        and
  \item repeated measurements establish how a specified workload
        behaved.
\end{enumerate}

The primary contrast holds the processor, socket, cores, executable,
model, input and operating conditions constant while changing verified
page placement between local HBM and local DDR5 in Flat mode.  This is
closer to a causal HBM comparison than either a Max 9480--Platinum 8592+
comparison or a Flat--Cache comparison.  The former changes processor
architecture; the latter changes the memory-management policy as well
as the tier.

\subsection{Experimental sequence}

The work should advance only after the preceding stage is stable.

\begin{figure}[tbp]
\centering
\begin{tikzpicture}[
  font=\small,
  stage/.style={draw=black!60, fill=neutralcol, rounded corners=2pt,
                text width=3.0cm, minimum height=1.15cm, align=center},
  primary/.style={stage, draw=accent, very thick, fill=accent!6},
  optional/.style={draw=black!45, dashed, fill=white, rounded corners=2pt,
                   text width=3.5cm, minimum height=0.95cm, align=center},
  flow/.style={-{Stealth}, thick, black!60}
]
  \node[stage] (s0) at (0,0)
    {\textbf{0. Manifest}\\topology, modes, software};
  \node[stage] (s1) at (3.7,0)
    {\textbf{1. Memory baseline}\\bandwidth, latency, locality};
  \node[primary] (s2) at (7.4,0)
    {\textbf{2. Primary contrast}\\local HBM vs.\ local DDR5};
  \node[stage] (s3) at (11.1,0)
    {\textbf{3. Boundaries}\\capacity and topology};
  \node[optional] (s4) at (11.1,-2.0)
    {\textbf{4. Conditional extension}\\energy, second runtime or code change};

  \draw[flow] (s0) -- (s1);
  \draw[flow] (s1) -- (s2);
  \draw[flow] (s2) -- (s3);
  \draw[flow] (s3) -- (s4);
\end{tikzpicture}
\caption{Experimental dependency order.  The primary HBM--DDR result
does not depend on a later implementation contribution.  Author's
design.}
\label{fig:experimental-sequence}
\end{figure}

\begin{description}
  \item[Stage 0: system manifest.] Record the items in
  Table~\ref{tab:platform-record}, identify HBM and DDR NUMA nodes, and
  verify whether Flat/Cache and Quadrant/SNC4 settings are available.

  \item[Stage 1: memory characterization.] Measure local and remote read,
  copy and triad bandwidth; unloaded and loaded latency; core-count
  scaling; and concurrent HBM--DDR traffic.  Establish saturation and
  run-to-run variability before fitting a performance model.

  \item[Stage 2: primary inference contrast.] Freeze one runtime commit,
  compiler and build.  Separate prefill from decode and compare verified
  local HBM with local DDR5 on one socket.  This stage is the minimum
  complete empirical study.

  \item[Stage 3: capacity and topology.] Move the live set through the
  local HBM boundary, then compare explicit placement, two-socket
  execution and SNC4 where available.  Cache mode is included only if a
  controlled reboot/configuration can be arranged.

  \item[Stage 4: conditional extension.] Add energy, a second runtime or
  a software change only after the baseline identifies a repeatable
  effect and the remaining time supports a proper evaluation.
\end{description}

\subsection{Selecting factors without a Cartesian explosion}

The experiment matrix should be chosen from hypotheses and pilot data,
not by taking every combination of model, precision, length, batch,
core count and topology.

\begin{table}[tbp]
\centering
\small
\caption{Factors and the role of each in the experimental design.}
\label{tab:experimental-factors}
\begin{tabularx}{\textwidth}{@{}lYY@{}}
\toprule
\textbf{Factor} & \textbf{Primary levels} & \textbf{Purpose} \\
\midrule
Memory placement
 & Local HBM and local DDR5; remote and mixed placement in a later stage
 & Isolate tier bandwidth, then quantify locality and spill effects. \\
Working-set regime
 & Comfortably below, near, and above usable local HBM capacity
 & Locate a capacity boundary using measured resident memory rather than
   parameter count alone. \\
Precision
 & BF16 reference and at least one practical quantized weight format;
   KV precision recorded independently
 & Test whether reduced traffic, packing and kernel coverage change the
   HBM benefit. \\
Workload shape
 & Short and long prompt; short and sustained generation; low and
   moderate batch/concurrency
 & Separate prefill, weight-streaming decode and KV-dominant regimes. \\
Core/topology
 & Partial socket, full local socket, two local replicas or an explicit
   shard; SNC4 only if available
 & Identify saturation and distinguish additional compute from remote
   traffic. \\
Cache state
 & Cold-start/model-load and warmed steady state reported separately
 & Prevent file I/O or page-cache state from contaminating inference
   measurements. \\
\bottomrule
\end{tabularx}
\end{table}

Models should be selected by memory regime and architecture rather than
popularity alone.  The design should include a model/format that fits
comfortably in local HBM and one configuration close to the usable
capacity boundary.  At least one grouped-query model is useful because
its KV growth differs from full multi-head attention.  If the boundary
cannot be crossed with available models, context, concurrency or a
controlled allocation can be used to change the live set, provided the
mechanism does not alter the computation being compared.

Pilot runs determine the final numeric levels.  A sparse matrix of
representative points followed by focused sweeps around observed knees
is more informative than a large under-replicated grid.

\subsection{Runtime and correctness controls}

One runtime should be selected as the primary implementation and frozen
at a named commit.  Selection criteria are support for the required
model formats, a measurable prefill/decode boundary, controllable thread
affinity, transparent build reproducibility and a viable way to direct
or verify memory placement.  A second runtime is useful for external
validity but should not multiply every factor in the main experiment.

Performance comparisons are meaningful only when the computation is
equivalent.  Record the model revision and tokenizer, use fixed prompts
and seeds where supported, and validate token sequences or task-level
quality when precision changes.  If different kernels or quantization
formats change model quality, report that difference rather than
treating them as interchangeable.

\subsection{Metrics and explanatory measurements}

The primary metrics are:

\begin{itemize}
  \item prefill-dominated time to first output token (TTFT), with the
        timed boundary declared;
  \item median, p95 and p99 inter-token latency for decode;
  \item output tokens/s and requests/s at the declared concurrency; and
  \item peak and steady resident memory by object or tier where the
        runtime exposes it.
\end{itemize}

E2E latency should be reported with an explicit boundary.  Model loading
and tokenization are measured separately unless the experiment is
deliberately end-to-end.  Throughput and reciprocal latency from the
same fixed synchronous batch are two presentations of one observation,
not independent confirmation.

To explain a result, collect the smallest counter set that can test the
mechanism: HBM and DDR traffic, UPI or remote traffic, achieved
frequency, cache misses and relevant matrix-engine activity.  Counter
names and availability are processor- and permission-specific.  A
counter should be omitted rather than assigned a meaning that its scope
does not support.

\subsection{Run protocol and statistical analysis}

Each reported comparison should follow the same protocol:

\begin{enumerate}
  \item reserve an exclusive node and archive the system manifest;
  \item pin physical cores and memory explicitly, with SMT treated as a
        separate factor;
  \item verify page residency after model initialization and during the
        run, not only after allocation;
  \item separate model loading and warm-up from steady-state timing,
        unless they are part of the declared metric;
  \item counterbalance or randomize HBM and DDR run order to reduce
        thermal and temporal bias;
  \item retain every raw sample, failure and outlier flag together with
        the command and environment;
  \item determine replication from pilot variability, using paired runs
        on the same node where possible; and
  \item repeat a subset on more than one HBM node before generalizing to
        the partition.
\end{enumerate}

Report absolute values and paired effect sizes, not only normalized bar
charts.  For primary comparisons, provide the median and interquartile
range together with a 95\% confidence interval appropriate to the
sampling design.  Tail percentiles require enough generated tokens or
requests to estimate them; p99 from a handful of observations is not a
tail-latency result.

\subsection{Energy, if the boundary can be defended}

RAPL estimates processor-domain energy; XClarity/INA226 and a rack or
node power meter may cover different hardware.  Before reporting
joules/token, validate access, units, update rate, counter wraparound and
the components included in each source.  Integrate energy over the same
timed region as performance.  Gross joules/request and
joules/output-token should be reported; if idle energy is subtracted,
also report the gross value and the subtraction procedure.  Facility
cooling claims require a measured facility boundary and are outside a
node-level study.

\subsection{Threats to validity}

\begin{table}[tbp]
\centering
\small
\caption{Main validity threats and design responses.}
\label{tab:validity}
\begin{tabularx}{\textwidth}{@{}lYY@{}}
\toprule
\textbf{Validity} & \textbf{Threat} & \textbf{Response} \\
\midrule
Internal
 & Frequency drift, first touch, page migration, remote access, warm
   cache state or shared-node interference is mistaken for an HBM effect.
 & Exclusive nodes, paired order, explicit binding, residency checks and
   frequency/traffic measurements. \\
Construct
 & STREAM is treated as LLM bandwidth; file size as live memory; average
   TPOT as tail latency; RAPL as whole-node energy.
 & Define each metric and boundary, measure the live set, and use
   counters only to test a stated mechanism. \\
External
 & One model, precision, node or runtime is generalized to CPU inference
   as a whole.
 & Cover distinct memory regimes, repeat on multiple nodes, and state the
   platform/runtime boundary of every conclusion. \\
Statistical conclusion
 & Few repetitions, selected best runs or unreported failures inflate an
   effect.
 & Retain raw samples, predeclare primary comparisons and report
   uncertainty with absolute values. \\
Reproducibility
 & Modules, firmware, models and scheduler inventory change over time.
 & Archive immutable source/model revisions, build metadata, commands
   and dated system manifests. \\
\bottomrule
\end{tabularx}
\end{table}

\FloatBarrier

\subsection{When an implementation contribution is justified}

The empirical study is complete without a new placement algorithm.  A
software contribution should follow only if the baseline reveals a
stable inefficiency, such as remote pages under an ostensibly local
policy, an avoidable capacity spill, or poor use of the second socket.
The smallest defensible intervention is preferable: separate
HBM/DDR allocators, object-level placement, one worker per locality
domain, or a topology-aware request router.

Any intervention must be compared with strong static baselines,
including all-local placement when the model fits, weights-first and
KV-first tiering, independent per-socket replicas, and hardware Cache
mode if available.  Its evaluation must include placement overhead,
remote traffic, memory use and output quality.  Novelty should be
claimed only after a current literature search and an evaluation show
that the measured advantage is not already provided by an existing
runtime.  It must also distinguish an implemented contribution from the
unevaluated NUMA-aware placement proposed by Na et al.
~\cite{na2024llmcpu}.
